% !TeX root = ../papers/nc-p3.tex
\section{長さ四の全Artin切断と核側の反対対称性}
\label{sec:finite-artin-kernel-symmetry}
\begin{theorem}\label{thm:finite-artin-kernel-symmetry}
定理\ref{thm:finite-tetrahedral-kernel-symmetry}で $\Sigma$ は長さ4であればよく、
非被約でも奇位置の反対対称性が成立する。従って四辺が有限、反対階数8で、
接次元1以下の局所成分が一つでもある候補は、中央支持3の点から正則性が決まるか、
少なくとも片側の反対対称性を持つ。いずれも全次数の標準正則性が有限に決まる。
\end{theorem}
\begin{proof}[外積の環と線型な分母]
$B=H^0(\mathcal O_\Sigma)$ とし、定理\ref{thm:finite-marked-artin-secants}の完全跡表示を使う。
偶座標を $a\in B$ と書くと、その対称行列は $\tau(a u v)$ ($u,v\in V=\ker\tau$) である。
点環上で奇位置の射影座標 $l,r$ は $V$ に値を持ち、$al=\alpha 1$ があるスカラー $\alpha$ で成り立つ。
$V$ の非零元の乗法階数は2以上である。階数1ならその元の張る直線がイデアルとなるが、
完全跡の核は非零イデアルを含まない。

非対称を仮定し、前定理と同じく $k=l\times r$ の成分イデアルを $(\eta)$ とし、
$S=D/\operatorname{ann}(\eta)$ と置く。局所生成元は直積環上の一元にまとめられる。
前定理の外積の恒等式から $Z=\operatorname{Spec}S\subset\Sigma$ であり、$S=B/I$ でもある。
$Z$ 上で二つの偶行列の全列は同じ直線に入り、対称性と単元成分により両行列は単元倍で一致する。
$\Sigma$ の $\mathbb P(B)$ での普遍点 $a$ は、完全跡により評価指標に双対な直線である。
実際指標 $\chi$ の双対元は $b a=\chi(b)a$ を満たし、制限形式は
$\chi(u)\chi(v)$ となる。この写像が完全な $B$ の座標系による閉埋込みで、
長さ4の切断に含まれることから切断全体に一致する。
従って $Z$ 上で $a$ の全乗法行列 $m_a$ も階数1である。
$k=\eta (u_1,u_2,u_3)$ の $u_i\in S$ は、単元倍を $\eta$ に吸収すれば
$V$ の三基底の $S$ への像である。$\dim S=n$ は3または4、印はその台にない。
$S$ は $D$ の完全跡に $\eta$ を掛けて得る完全積対合によりGorensteinである。

位置3の四座標の $S$ への像を $W'$ とする。$h^{-1}W^2=U^2$ として
$h^{-1}\eta s$ の一意な位置1三次持上げを $\Phi(s)$ と書く。
$\Phi$ は単射で、$S$ の完全跡 $\tau_S$ により $\Phi(s)=\tau_S(Cs)$ と表せる。
ここで $C$ は位置1の三次同次式を係数とする $S$ の元である。
独立変数での $l$ の $S$ への射影を $L$ とする。前定理と同じ外積に $l$ を掛けた恒等式より
$\tau_S(LC b)=0$ ($b\in W'$) である。一般体上で $C$ は単元である。
そうでなければ、ある局所因子の定数socle元を $\Phi$ が消す。$L$ も一般体上で単元である。
よって $W'\ne S$ で、その次元は閉埋込みより2以上である。

まず $W'$ が余次元1なら $LC=\lambda c$、$c$ はその直交補を張る定数単元、$\lambda$ は四次式である。
$S$ が曲線状の局所因子の積なら、$L^{-1}$ の最小共通分母は $N=\det m_L$（定数倍を除く）。
実際各因子 $k[e]/e^d$ で $L=L_0+eL_1+\cdots$ と書くと、$L_0$ に沿う極は正確に $d$ 次である。
$d\ge2$ では方向射影 $\Sigma\to\mathbb P(V^*)$ の閉埋込みにより $L_1$ は $L_0$ に比例せず、
逆数展開の $e^{d-1}$ 係数の分子は法 $L_0$ で $(-L_1)^{d-1}$ となる。
この閉埋込みは異なる二点と全接方向を $V$ が分離することによる
（任意の長さ3の部分が三座標を張る）。異なる局所因子の $L_0$ は比例しない。
従って $C=H c\,N L^{-1}$、$\deg H=4-n$ である。$NL^{-1}$ の乗法行列は
$\operatorname{adj}m_L$ であり、全係数を実際の点環に戻すと
$H\operatorname{adj}m_L$ の成分は $(\eta)$ に入る（$C$ の係数関数は $\Phi(S)$ を張る）。
\end{proof}

\begin{proof}[乗法の階数と印での長さ]
$n=3$ のGorenstein環は曲線状の因子の積である（局所埋込次元2なら長さ3では平方零となる）。
$I$ は長さ1のイデアルで、$B$ は $I$ に指標 $\chi$ で作用する。
印以外で $H\chi(a)\in(\eta)$ を示す。まず $(\eta)$ を法とし、さらに $H$ の零化イデアルで割る。
この局所環上で $\operatorname{adj}m_L=0$、すなわち $S$ 上の乗法の全二次小行列式は零。
一方 $l$ の剰余は非零で、$B$ での乗法には単元な二次小行列式がある。
$B\to S$ の核は一次元なので、$S$ 上での像は階数1の直和因子である。
$B$ 上の $lB$ はこの直線の逆像全体（階数2の直和因子）である。
実際その逆像に入り、剰余で二つの独立列を持つためである。特に $I$ を含む。
$a(lB)=\alpha B$ は二元で生成されるため $\alpha=0$。非零なら $\alpha B$ は基底環の
非零な主イデアルの四コピーで、Nakayamaにより四生成元を要する。
従って $aI=0$、すなわち $\chi(a)=0$ であり、元の環に戻して所望の包含を得る。
印 $k[t]/t^m$ で $m>2$ なら $l/t$ を法 $t^{m-2}$ で同じ議論に用いる。
$H\operatorname{adj}m_L$ の次数2の因子を $t^2$ で割れ、$a(l/t)$ もスカラーである。
よって $H\chi(a)$ の支持長は $3+2$ 以下である。非零な隣接語の下界6と純関係の不在に反する。

$n=4$, $S=B$ が曲線状の因子の積なら $H$ は定数である。印以外で $(\eta)$ を法とすると
$\operatorname{adj}m_l=0$ かつ $m_l$ に単元な二次小行列式がある。従って像は階数2の
直和因子で、先と同じ生成元数から $\alpha=0\bmod(\eta)$ である。
元の局所環でも $lB$ の二列を基底の二元に取れる。$m_a$ はそれらを $\alpha$ の倍に移すので、
三次小行列式は $\alpha$ と $m_a$ の二次小行列式のイデアルとの積に入る。
列の基底変更を行ってその列で展開すればよい。
$Z$ 上で $m_a$ が階数1であったため、この積は零である。
従って印以外で $\operatorname{adj}m_a=0$。印で $m>3$ なら $l/t$ と法 $t^{m-3}$ での
同じ階数2の像を使い、$a$ がそれを消すのでやはり $\operatorname{adj}m_a=0\bmod t^{m-3}$。
しかし $N(a)a^{-1}$ の四座標は独立な三次式である。単元の反転が射影単元開集合を
全体へ移すことから従う。その四式が長さ3以下のイデアルに入るのは不可能である。

長さ4の残る非曲線状Gorenstein環 $B$ は局所で Hilbert 関数 $(1,2,1)$ を持つ。
実際埋込次元3なら平方零である。極大イデアルの二乗はsocle $k\omega$、三乗は零で、
一次の二座標上の乗法は非退化二次形式である。
$V\to B/k\omega$ は同型なので、$l=s+v$（$v$ は根基に値を持つ）に対し
逆数の最小分母は $s^3$ で、分子は $s^2-sv+v^2$。
二次形式 $v^2/\omega$ が法 $s$ で零にならないことを用いた。
従って $C=H c(s^2-sv+v^2)$、$H$ は非零一次式である。
剰余指標を $\chi$ とすると印以外では $H\chi(a)\in(\eta)$ である。
実際法 $(\eta)$ で $s$ が単元なら $H=0$。そうでなければ $v$ の一次二座標の一つが単元で、
$Hsv=0$ を根基の二乗を法として読めば $Hs=0$。
$al=\alpha 1$ の根基の一次部分に $H$ を掛けると $H\chi(a)v=0$ であり結論を得る。
印で $m>2$ なら次数2の分子を $t^2$ で割り同じ議論を行う。
従って隣接語 $H\chi(a)$ の支持長は $4+2$ 以下。6未満なら不可能で、6なら
印以外で長さ4を全て使い、その支持は $Z$ を含む。$Z=\Sigma$ の偶座標は全四空間を張るが、
補題\ref{lem:finite-adjacent-word-support}の長さ6の語はどちらの座標でも平面に入る。これも矛盾する。
\end{proof}

\begin{proof}[長さ四の共線スキーム]
最後に $n=4$, $\dim W'=2$ とする。$Z$ は位置3の直線における長さ4の因子である。
無限遠を外し $S=k[t]/P(t)$ とし、$P$ は四次モニックとする。両隣の偶座標はそれぞれ
$Z$ に全射する完全一次系なので、個別の自明化と基底で両方を
$e(t)=(1,t,t^2,t^3)^t$ と書ける。$t e=T e\bmod P$ と定数行列 $T$ を取る。
直線上の四テンソルを、残る一位置の任意の縮約ごとに $M(t)=M_0+tM_1$ として両偶位置の間で読む。
隣接条件 $e(t)^t M(t)=M(t)e(t)=0\bmod P$ と四基底の独立性より
$M_0=-T^t M_1=-M_1T$。$P$ の根 $d$ で $Te(d)=d e(d)$ となる。
従って $e(d)^t M_1$ は $T$ の固有値 $d$ の左固有ベクトルまたは零で、
この固有空間は巡回行列のため重根でも一次元である。
よって前位置を $e(d)$ に、中央位置を直線の $t\ne d$ に制限すれば、
残る位置の全縮約の後位置での像は固定直線に入る。
この二文字縮約は零ではない。零なら $e(d)$ が点辺で中央直線の全点を後続に持ち、
辺の有限性に反する。従って後位置の関係空間に純関係があり矛盾する。

これで長さ4の全切断が尽きる。反対階数8の系では各側で共通切断が自動であり、
有限長3以下と正次元の切断では先の各割線・母線の定理が切断の偶座標の対称性を与える。
長さ4では今回の奇側の結論を使えばよい。片側対称になった場合の全次数判定は
系\ref{cor:all-defective-one-sym-finite}を適用する。具体的には対称核に階数2以下の元があれば認証され、
それ以外は両対称窓の有限判定で決まる。
\end{proof}
